\begin{textsample}{POS Dim 3 – ro – Score 17.00 – ro/ro\_inf\_33.txt}  \label{ex:f3_pos_129}
Direitos de aprendizagem e desenvolvimento deste campo, articulados aos demais campos - Conviver com \textbf{crianças} e \textbf{adultos} compartilhando sua língua materna em situações comunicativas cotidianas, constituindo modos de pensar, imaginar, sentir, narrar, dialogar e conhecer. - \textbf{Brincar} com parlendas, trava-línguas, adivinhas, memória, rodas, \textbf{brincadeiras} cantadas, jogos e textos de imagens, escritos e outros, ampliando o repertório das manifestações culturais da tradição local e de outras culturas, enriquecendo sua linguagem oral, corporal, musical, \textbf{dramática}, escrita, dentre outras. - Explorar \textbf{gestos}, expressões, \textbf{sons} da língua, rimas, imagens, textos escritos, além dos sentidos das palavras, nas poesias, parlendas, canções e nos enredos de histórias, apropriando -se desses elementos para criar novas falas, enredos, histórias e escritas convencionais ou não. - Participar de rodas de conversa, de relatos de experiências, de contação e leitura de histórias e poesias, de construção de narrativas, da elaboração, descrição e representação de papéis no faz de conta, da exploração de materiais impressos e de variedades linguísticas, construindo diversas formas de organizar o pensamento. - Expressar sentimentos, ideias, percepções, \textbf{desejos}, necessidades, pontos de vista, informações, dúvidas e descobertas utilizando múltiplas linguagens, entendendo e considerando o que é comunicado pelas demais \textbf{crianças} e \textbf{adultos}. - Conhecer -se e reconhecer suas preferências por pessoas, \textbf{brincadeiras}, lugares, histórias, autores, gêneros linguísticos, e seu interesse em produzir com a linguagem verbal. \textbf{Escuta}, Fala, Pensamento e Imaginação Contextos de Aprendizagem e Desenvolvimento Para propor intencionalidades educativas significativas, docentes da primeira infância podem : Pensando nos \textbf{Bebês} - Propor situações de conversa com os \textbf{bebês} em ambiente tranquilo e \textbf{lúdico}. - Favorecer participação em jogos rítmicos em que os \textbf{bebês} são animados a \textbf{imitar} \textbf{sons} variados, ou em jogos de nomeação em que a \textbf{professora} aponta para algo, propõe a questão : “ O que é isso? ”, e \textbf{apoia} o \textbf{bebê} a responder. - Oportunizar \textbf{brincadeiras} com outros \textbf{bebês}, com ou sem objetos, expressando -se, corporal e/ou verbalmente. - Possibilitar que repitam acalantos, cantigas de roda, poesias e parlendas, explorando o ritmo, a sonoridade e a conotação das palavras. - Oferecer a \textbf{escuta} de histórias, contos de repetição, poemas e \textbf{imitação} das variações de entonação e de \textbf{gestos} realizados pela \textbf{professora} ao ler ou cantar. - Oportunizar situações de \textbf{brincar} de traçar marcas gráficas em cartolinas ou outro suporte, usando os \textbf{dedos} ou pincéis, além de \textbf{tintas} caseiras ou não.

% matched lemmas: adulto, apoiar, bebê, brincadeira, brincar, criança, dedo, desejo, dramático, escuta, gesto, imitar, imitação, lúdico, professora, som, tinta
\end{textsample}
